\begin{thebibliography}{99}
\small
\raggedright
\setlength{\itemsep}{2pt}
\bibitem{erdos1948} P.~Erd\H{o}s,
  \href{https://users.renyi.hu/~p_erdos/1948-04.pdf}{\emph{On arithmetical
  properties of Lambert series}},
  J. Indian Math. Soc. (N.S.) \textbf{12} (1948), 63--66.
\bibitem{erdos1988} P.~Erd\H{o}s, \emph{On the irrationality of certain series:
  problems and results}, in A.~Baker (ed.), \emph{New Advances in Transcendence
  Theory}, Cambridge UP, 1988, pp.~102--109,
  doi:\href{https://doi.org/10.1017/CBO9780511897184.009}{10.1017/CBO9780511897184.009}.
\bibitem{borwein1991} P.~B. Borwein, \emph{On the irrationality of
  $\sum1/(q^{n}+r)$},
  J. Number Theory \textbf{37} (1991), no.~3, 253--259,
  doi:\href{https://doi.org/10.1016/S0022-314X(05)80041-1}{10.1016/S0022-314X(05)80041-1}.
\bibitem{borwein1992} P.~B. Borwein, \emph{On the irrationality of certain
  series}, Math. Proc. Cambridge Philos. Soc. \textbf{112} (1992), no.~1,
  141--146,
  doi:\href{https://doi.org/10.1017/S030500410007081X}{10.1017/S030500410007081X}.
  Van Assche cites Lemma~2 for the neighbouring little-$q$-Legendre
  evaluation used above.
\bibitem{az1998} T.~Amdeberhan and D.~Zeilberger, \emph{$q$-Ap\'ery
  irrationality proofs by $q$-WZ pairs}, Adv. Appl. Math. \textbf{20} (1998),
  no.~2, 275--283,
  doi:\href{https://doi.org/10.1006/aama.1997.0565}{10.1006/aama.1997.0565};
  arXiv:\href{https://arxiv.org/abs/math/9804122v1}{math/9804122v1}.
  Page references are to arXiv:math/9804122v1.
\bibitem{bv1994} P.~Bundschuh and K.~V\"a\"an\"anen,
  \href{https://numdam.org/item/CM_1994__91_2_175_0.pdf}{\emph{Arithmetical
  investigations of a certain infinite product}},
  Compositio Math. \textbf{91} (1994), no.~2, 175--199.
\bibitem{zudilin2002} W.~Zudilin, \emph{Remarks on irrationality of
  \(q\)-harmonic series}, Manuscripta Math. \textbf{107} (2002), no.~4,
  463--477,
  doi:\href{https://doi.org/10.1007/s002290200249}{10.1007/s002290200249}.
\bibitem{zudilin2004} W.~Zudilin,
  \href{https://geodesic.mathdoc.fr/articles/10.4064/aa111-2-4/}{\emph{Heine's
  basic transform and a permutation group for \(q\)-harmonic series}},
  Acta Arith. \textbf{111} (2004), no.~2, 153--164,
  doi:\href{https://doi.org/10.4064/aa111-2-4}{10.4064/aa111-2-4}.
  Page references are to the printed journal pages.
\bibitem{zudilin2016} W.~Zudilin,
  \href{https://arxiv.org/abs/1601.02688v2}{\emph{On the irrationality of
  generalized \(q\)-logarithm}}, arXiv:1601.02688; Res. Number Theory
  \textbf{2} (2016), Art.~15,
  doi:\href{https://doi.org/10.1007/s40993-016-0042-x}{10.1007/s40993-016-0042-x}.
  Page references are to arXiv:1601.02688v2.
  The remark that the results extend to non-integer \(p=r/s\), \(|p|>1\), under
  an assumption \(\log|r|>c\log|s|\) for a computable \(c>0\), is in Section~2,
  p.~4, in the paragraph beginning ``Finally, we remark''; no value of \(c\) is
  computed there, and the remark is made for the generalized \(q\)-logarithm of
  that paper.
\bibitem{zudilin2017det} W.~Zudilin, \emph{A determinantal approach to
  irrationality}, Constr. Approx. \textbf{45} (2017), no.~2, 301--310,
  doi:\href{https://doi.org/10.1007/s00365-016-9333-7}{10.1007/s00365-016-9333-7};
  arXiv:\href{https://arxiv.org/abs/1507.05697v1}{1507.05697v1}.
  Page and equation references are to arXiv:1507.05697v1.
\bibitem{bundschuhzudilin2008} P.~Bundschuh and W.~Zudilin, \emph{Rational
  approximations to a \(q\)-analogue of \(\pi\) and some other \(q\)-series},
  in H.~P. Schlickewei, K.~Schmidt and R.~F. Tichy (eds.), \emph{Diophantine
  Approximation}, Dev. Math. \textbf{16}, Springer, 2008, pp.~123--139,
  doi:\href{https://doi.org/10.1007/978-3-211-74280-8_6}{10.1007/978-3-211-74280-8\_6}.
\bibitem{krvz2009} C.~Krattenthaler, I.~Rochev, K.~V\"a\"an\"anen and
  W.~Zudilin, \emph{On the non-quadraticity of values of the
  \(q\)-exponential function and related \(q\)-series}, Acta Arith.
  \textbf{136} (2009), no.~3, 243--269,
  doi:\href{https://doi.org/10.4064/aa136-3-4}{10.4064/aa136-3-4};
  arXiv:\href{https://arxiv.org/abs/0812.2921v1}{0812.2921v1}.
  Page references are to arXiv:0812.2921v1.
\bibitem{duverney1996} D.~Duverney,
  \href{https://numdam.org/item/JTNB_1996__8_1_173_0.pdf}{\emph{\`A propos de la
  s\'erie $\sum_{n\ge1}x^{n}/(q^{n}-1)$}},
  J. Th\'eor. Nombres Bordeaux \textbf{8} (1996), no.~1, 173--181.
  Th\'eor\`eme~2 on p.~174 gives the rational-base region
  $\log|s|/\log|r|<\frac13(1-3/\pi^{2})=0.2320\ldots$ for this series;
  Th\'eor\`eme~1, for a general numerator, is weaker.
\bibitem{matalaaho2006} T.~Matala-aho, K.~V\"a\"an\"anen and W.~Zudilin,
  \href{https://doi.org/10.1090/S0025-5718-05-01812-0}{\emph{New irrationality
  measures for \(q\)-logarithms}}, Math. Comp. \textbf{75} (2006), no.~254,
  879--889, doi:10.1090/S0025-5718-05-01812-0.  The hypothesis
  \(p=1/q\in\Z\mathbin{\backslash}\{0,\pm1\}\) is carried in the abstract on p.~879 and in
  both theorem statements on p.~880, where the authors also record that their
  methods do not sharpen the \(q\)-harmonic case of~\cite{zudilin2004}.
\bibitem{adamczewskibell2013} B.~Adamczewski and J.~P. Bell,
  \href{https://arxiv.org/abs/1303.2019v1}{\emph{A problem about Mahler
  functions}}, Ann. Sc. Norm. Super. Pisa Cl. Sci. \textbf{17} (2017), no.~4,
  1301--1355; arXiv:\href{https://arxiv.org/abs/1303.2019v1}{1303.2019v1}, 2013.
  Theorem~1.1 on p.~6 of arXiv:1303.2019v1: over a field of characteristic
  zero, a power series is both \(k\)- and \(\ell\)-Mahler for multiplicatively
  independent \(k,\ell\) if and only if it is a rational function.  Bell and
  Smertnig cite it as Theorem~1.3 of the journal version
  \cite[Thm.~2.5, p.~5]{bellsmertnig2026}.
\bibitem{rhinviola1996} G.~Rhin and C.~Viola,
  \href{https://geodesic.mathdoc.fr/articles/10.4064/aa-77-1-23-56/}
  {\emph{On a permutation group related to $\zeta(2)$}},
  Acta Arith. \textbf{77} (1996), no.~1, 23--56,
  doi:10.4064/aa-77-1-23-56.
\bibitem{vanassche2001} W.~Van Assche,
  \href{https://arxiv.org/abs/math/0101187v1}{\emph{Little \(q\)-Legendre
  polynomials and irrationality of certain Lambert series}},
  Ramanujan J. \textbf{5} (2001), no.~3, 295--310,
  doi:\href{https://doi.org/10.1023/A:1012930828917}{10.1023/A:1012930828917}.
  Page references are to arXiv:math/0101187v1.
\bibitem{coussementsmet2007} J.~Coussement and C.~Smet, \emph{Irrationality
  proof of certain Lambert series using little \(q\)-Jacobi polynomials},
  J. Comput. Appl. Math. \textbf{233} (2009), no.~3, 680--690,
  doi:\href{https://doi.org/10.1016/j.cam.2009.02.036}{10.1016/j.cam.2009.02.036};
  arXiv:\href{https://arxiv.org/abs/math/0701345v1}{math/0701345v1}.
  Page references are to arXiv:math/0701345v1.
\bibitem{koizumiyokoi2026} J.~Koizumi and A.~Yokoi, \emph{Ap\'ery-type
  approximations and irrationality measures for certain \(q\)-series},
  arXiv:\href{https://arxiv.org/abs/2608.26918v1}{2608.26918v1},
  27 August 2026.
\bibitem{vandehey2013} J.~Vandehey,
  \href{https://arxiv.org/abs/1206.0340v1}{\emph{On an incomplete argument of
  Erd\H{o}s on the irrationality of Lambert series}}, Integers
  \textbf{13} (2013), Paper A58. Page references are to
  arXiv:1206.0340v1 (2012).
\bibitem{lucatachiya2017} F.~Luca and Y.~Tachiya,
  \href{https://www.kurims.kyoto-u.ac.jp/~kyodo/kokyuroku/contents/pdf/2014-14.pdf}
  {\emph{Linear independence results for the values of divisor functions
  series}}, RIMS K\^oky\^uroku No.~2014 (2017), 138--150.
  Theorem~A on p.~139 restates the periodic-coefficient irrationality theorem;
  Example~1 on p.~140 gives the divisor-function specialization.
\bibitem{duverneytachiya2019} D.~Duverney and Y.~Tachiya, \emph{Refinement of
  the Chowla--Erd\H{o}s method and linear independence of certain Lambert
  series}, Forum Math. \textbf{31} (2019), no.~6, 1557--1566,
  doi:\href{https://doi.org/10.1515/forum-2018-0299}{10.1515/forum-2018-0299}.
  Page references are to the
  \href{https://danielduverney.fr/documents/theorie-des-nombres/DuverneyTachiya190522.pdf}{authors'
  version}.
\bibitem{rivin2026} I.~Rivin,
\href{https://arxiv.org/abs/2604.25151v1}{\emph{Zero Coefficients of Rational
  Power Series and Rational Lambert Series}}, arXiv:2604.25151v1,
  28 April 2026.
  Theorem~1.1 is on p.~2 and proved on pp.~6--7; the periodic-coefficient
  Corollary~6.4 is on p.~9.
\bibitem{kovactao2024} V.~Kova\v{c} and T.~Tao,
  \href{https://arxiv.org/abs/2406.17593v4}{\emph{On several irrationality
  problems for Ahmes series}}, Acta Math. Hungar. \textbf{175} (2025), no.~2,
  572--608,
  doi:\href{https://doi.org/10.1007/s10474-025-01528-0}{10.1007/s10474-025-01528-0};
  arXiv:\href{https://arxiv.org/abs/2406.17593v4}{2406.17593v4}.
  Page references are to arXiv:2406.17593v4.
\bibitem{erdosproblems} T.~F. Bloom,
  \href{https://www.erdosproblems.com/1049}{\emph{Erd\H{o}s Problem \#1049}},
  \texttt{erdosproblems.com/1049}. Historical snapshot cited in the supplied
  manuscript: accessed 28 July 2026, displaying ``last edited
  28 September 2025''.
\bibitem{bellsmertnig2026} J.~Bell and D.~Smertnig,
  \href{https://arxiv.org/abs/2603.23456v1}{\emph{Mahler series with
  multiplicative coefficient sequences}}, arXiv:2603.23456v1,
  24 March 2026.  Theorem~1.3 is on pp.~2--3; its stated consequences on
  p.~3 include that the divisor and totient generating series are not
  $k$-Mahler for any $k\ge2$.
\bibitem{dlmf} F.~W.~J. Olver et al.\ (eds.), \emph{NIST Digital Library of
  Mathematical Functions},
  \href{https://dlmf.nist.gov/17.2.E37}{\texttt{dlmf.nist.gov/17.2.E37}},
  Eq.~17.2.37 in \S17.2(iii), accessed 15 September 2026.
\bibitem{formalconjectures1049} The Formal Conjectures Authors,
  \href{https://github.com/google-deepmind/formal-conjectures/blob/f776d2f2039351b00737ffcafb9d7d7666e1d9af/FormalConjectures/ErdosProblems/1049.lean}
  {\emph{FormalConjectures.ErdosProblems.\texttt{1049}}},
  Lean source at commit \texttt{f776d2f}, 2026, accessed 28 July 2026.  The
  irrationality declarations are unproved; the Lambert-series identity is
  proved.
\bibitem{postelmansvanassche2007} K.~Postelmans and W.~Van Assche,
  \href{https://arxiv.org/abs/math/0604312v1}{\emph{Irrationality of
  $\zeta_q(1)$ and $\zeta_q(2)$}}, J. Number Theory \textbf{126} (2007),
  no.~1, 119--154,
  doi:\href{https://doi.org/10.1016/j.jnt.2006.11.011}{10.1016/j.jnt.2006.11.011}.
  Page references are to arXiv:math/0604312v1 (2006).
\bibitem{krattenthaler2021} C.~Krattenthaler,
  \href{https://arxiv.org/abs/2103.03969v1}{\emph{A determinant identity for
  moments of orthogonal polynomials that implies Uvarov's formula for the
  orthogonal polynomials of rationally related densities}},
  arXiv:2103.03969v1 (2021). Page references are to this version.
\bibitem{wangzhu2016} Y.~Wang and B.-X.~Zhu,
  \href{https://arxiv.org/abs/1612.04114v1}{\emph{Log-convex and Stieltjes
  moment sequences}}, Adv. Appl. Math. \textbf{81} (2016), 115--127,
  doi:\href{https://doi.org/10.1016/j.aam.2016.06.008}{10.1016/j.aam.2016.06.008}.
  Page references are to arXiv:1612.04114v1.
\bibitem{stanley2016} R.~P. Stanley, \emph{Smith normal form in combinatorics},
  J. Combin. Theory Ser. A \textbf{144} (2016), 476--495,
  doi:\href{https://doi.org/10.1016/j.jcta.2016.06.013}{10.1016/j.jcta.2016.06.013};
  arXiv:\href{https://arxiv.org/abs/1602.00166v1}{1602.00166v1}.
  Page references are to arXiv:1602.00166v1.
\bibitem{duverney2011} D.~Duverney,
  \href{https://doi.org/10.1063/1.3630035}{\emph{Arithmetical functions and
  irrationality of Lambert series}}, AIP Conf. Proc. \textbf{1385} (2011),
  5--16, doi:10.1063/1.3630035. Theorem and page references are to the
  12-page author manuscript supplied with the research archive.
\bibitem{berg2007} C.~Berg, \emph{On powers of Stieltjes moment sequences, II},
  J. Comput. Appl. Math. \textbf{199} (2007), 23--38;
  arXiv:\href{https://arxiv.org/abs/math/0412340v1}{math/0412340v1}.
  Theorem references use the preprint; Theorem~5.1 treats factorial powers.
\bibitem{sw2024} A.~D. Sokal and J.~Walrad,
  \emph{Continued-fraction characterization of Stieltjes moment sequences
  with support in $[\xi,\infty)$},
  \href{https://arxiv.org/abs/2404.12131v1}{arXiv:2404.12131v1}, 2024.
  The classical Stieltjes criterion is recalled on pp.~1--2.
\bibitem{lrz2017} H.~Liang, J.~Remmel and S.~Zheng,
  \emph{Stieltjes moment sequences of polynomials},
  \href{https://arxiv.org/abs/1710.05795v1}{arXiv:1710.05795v1}, 2017.
\bibitem{psz2023} M.~P\'etr\'eolle, A.~D. Sokal and B.-X.~Zhu,
  \emph{Lattice paths and branched continued fractions: An infinite sequence
  of generalizations of the Stieltjes--Rogers and Thron--Rogers polynomials,
  with coefficientwise Hankel-total positivity}, Mem. Amer. Math. Soc.
  \textbf{291} (2023), no.~1450,
  doi:\href{https://doi.org/10.1090/memo/1450}{10.1090/memo/1450};
  \href{https://arxiv.org/abs/1807.03271v2}{arXiv:1807.03271v2}.
  Theorem and page references use that preprint version.
\bibitem{middeke2021} J.~Middeke, D.~J. Jeffrey and C.~Koutschan,
  \emph{Common Factors in Fraction-Free Matrix Decompositions},
  Math. Comput. Sci. \textbf{15} (2021), no.~4, 589--608,
  doi:\href{https://doi.org/10.1007/s11786-020-00495-9}{10.1007/s11786-020-00495-9};
  \href{https://arxiv.org/abs/2005.12380v1}{arXiv:2005.12380v1}.
  Section and theorem references use the preprint.
\bibitem{golubwelsch1969} G.~H. Golub and J.~H. Welsch,
  \emph{Calculation of Gauss Quadrature Rules}, Math. Comp. \textbf{23}
  (1969), no.~106, 221--230,
  doi:\href{https://doi.org/10.1090/S0025-5718-69-99647-1}{10.1090/S0025-5718-69-99647-1}.
\end{thebibliography}

\end{document}
